% from ega1-10-fr.tex
\oldpage[I]{213}
(\hyperref[I.5.5.9-fr]{5.5.9}) et
(\hyperref[I.5.5.10-fr]{5.5.10}) (en y remplaçant «~ouvert affine~» par
«~ouvert formel affine vérifiant la condition b) de
(\hyperref[I.10.6.3-fr]{10.6.3})~»).

Un raisonnement analogue montre aussi que tout schéma formel affine
\emph{noethérien} est séparé, ce qui justifie la terminologie.

\begin{proposition}[10.15.4]
\label{I.10.15.4-fr}
Soient $\mathfrak{S}$ un préschéma formel localement noethérien,
$\mathfrak{X}$, $\mathfrak{Y}$ deux $\mathfrak{S}$-préschémas formels
localement noethériens, tels que $\mathfrak{X}$ ou $\mathfrak{Y}$ soit de
type fini sur $\mathfrak{S}$ (de sorte que
$\mathfrak{X}\times_{\mathfrak{S}}\mathfrak{Y}$ est localement noethérien)
et que $\mathfrak{Y}$ soit séparé sur $\mathfrak{S}$. Soit
$f:\mathfrak{X}\to\mathfrak{Y}$ un $\mathfrak{S}$-morphisme~; alors le
morphisme graphe
$\Gamma_f=(1_{\mathfrak{X}},f)_{\mathfrak{S}}:\mathfrak{X}\to
\mathfrak{X}\times_{\mathfrak{S}}\mathfrak{Y}$ est une immersion fermée.
\end{proposition}

On peut supposer que $\mathfrak{S}$ est limite inductive d'une suite
$(S_n)$ de préschémas localement noethériens, $\mathfrak{X}$ (resp.
$\mathfrak{Y}$) limite inductive d'une suite $(X_n)$ (resp. $(Y_n)$) de
$S_n$-préschémas, $f$ limite inductive d'une suite de $S_n$-morphismes
$f_n:X_n\to Y_n$~; alors
$\mathfrak{X}\times_{\mathfrak{S}}\mathfrak{Y}$ est limite inductive de la
suite $(X_n\times_{S_n}Y_n)$ et $\Gamma_f$ de la suite
$(\Gamma_{f_n})$ (\hyperref[I.10.7.4-fr]{10.7.4})~; par hypothèse, $Y_0$ est
séparé sur $S_0$ (\hyperref[I.10.15.2-fr]{10.15.2}), donc l'espace
$\Gamma_{f_0}(X_0)$ est un sous-espace fermé de
$X_0\times_{S_0}Y_0$~; comme les espaces sous-jacents de
$\mathfrak{X}\times_{\mathfrak{S}}\mathfrak{Y}$ (resp.
$\Gamma_f(\mathfrak{X})$) et $X_0\times_{S_0}Y_0$ (resp.
$\Gamma_{f_0}(X_0)$) sont les mêmes, on voit déjà que
$\Gamma_f(\mathfrak{X})$ est un sous-espace \emph{fermé} de
$\mathfrak{X}\times_{\mathfrak{S}}\mathfrak{Y}$. Remarquons maintenant que
lorsque $(U,V)$ parcourt l'ensemble des couples formés d'un ouvert formel
affine noethérien $U$ (resp. $V$) de $\mathfrak{X}$ (resp.
$\mathfrak{Y}$) tels que $f(U)\subset V$, les ouverts
$U\times_{\mathfrak{S}}V$ forment un recouvrement de
$\Gamma_f(\mathfrak{X})$ dans
$\mathfrak{X}\times_{\mathfrak{S}}\mathfrak{Y}$, et si
$f_U:U\to V$ est la restriction de $f$ à $U$,
$\Gamma_{f_U}:U\to U\times_{\mathfrak{S}}V$ est la restriction de
$\Gamma_f$ à $U$. Si nous montrons que $\Gamma_{f_U}$ est une immersion
fermée, il en sera donc de même de $\Gamma_f$
(\hyperref[I.10.14.4-fr]{10.14.4}), autrement dit, on est ramené au cas où
$\mathfrak{S}=\operatorname{Spf}(A)$,
$\mathfrak{X}=\operatorname{Spf}(B)$,
$\mathfrak{Y}=\operatorname{Spf}(C)$ sont affines ($A$, $B$, $C$ adiques
noethériens), $f$ correspondant à un $A$-homomorphisme continu
$\varphi:C\to B$~; $\Gamma_f$ correspond alors à l'unique homomorphisme
continu $\omega:B\widehat{\otimes}_A C\to B$ qui, composé avec les
homomorphismes canoniques $B\to B\widehat{\otimes}_A C$ et
$C\to B\widehat{\otimes}_A C$, donne respectivement l'identité et
$\varphi$. Or, il est clair que $\omega$ est \emph{surjectif}, d'où notre
assertion.

\begin{corollary}[10.15.5]
\label{I.10.15.5-fr}
Soient $\mathfrak{S}$ un préschéma formel localement noethérien,
$\mathfrak{X}$ un $\mathfrak{S}$-préschéma de type fini~; pour que
$\mathfrak{X}$ soit séparé sur $\mathfrak{S}$, il faut et il suffit que le
morphisme diagonal
$\mathfrak{X}\to\mathfrak{X}\times_{\mathfrak{S}}\mathfrak{X}$ soit une
immersion fermée.
\end{corollary}

\begin{proposition}[10.15.6]
\label{I.10.15.6-fr}
Une immersion fermée $j:\mathfrak{Y}\to\mathfrak{X}$ de préschémas formels
localement noethériens est un morphisme séparé.
\end{proposition}

Avec les notations de (\hyperref[I.10.14.2-fr]{10.14.2}),
$j_0:Y_0\to X_0$ est une immersion fermée, donc un morphisme séparé, et il
suffit d'appliquer (\hyperref[I.10.15.2-fr]{10.15.2}).

\begin{proposition}[10.15.7]
\label{I.10.15.7-fr}
Soient $X$ un préschéma (usuel) localement noethérien, $X'$ une partie
fermée de $X$ et $\widehat{X}=X_{/X'}$. Pour que $\widehat{X}$ soit séparé,
il faut et il suffit que $\widehat{X}_{\mathrm{red}}$ le soit, et il suffit
que $X$ le soit.
\end{proposition}

En effet, avec les notations de (\hyperref[I.10.8.5-fr]{10.8.5}), pour que
$\widehat{X}$ soit séparé, il faut et il suffit que $X'_0$ le soit
(\hyperref[I.10.15.2-fr]{10.15.2}), et comme
$\widehat{X}_{\mathrm{red}}=(X'_0)_{\mathrm{red}}$, il est équivalent de
dire que $\widehat{X}_{\mathrm{red}}$ l'est
(\hyperref[I.5.5.1-fr]{5.5.1, (vi)}).
